\documentclass[11pt]{article}
\usepackage[a4paper,top=1.4cm,bottom=1.3cm,left=1.6cm,right=1.6cm]{geometry}
\usepackage[T1]{fontenc}
\usepackage{times}
\usepackage[scaled=0.88]{helvet}
\usepackage{microtype}
\usepackage{booktabs}
\usepackage{xcolor}
\usepackage{graphicx}
\usepackage{enumitem}
\usepackage{multicol}
\usepackage[hidelinks]{hyperref}
\pagestyle{empty}
\setlist{nosep,leftmargin=1.1em}
\setlength{\parindent}{0pt}
\setlength{\parskip}{3pt}
\setlength{\columnsep}{0.55cm}
\definecolor{ink}{HTML}{1B2A41}
\definecolor{acc}{HTML}{B4462B}
\newcommand{\onus}{\textsc{Onus}}
\newcommand{\hd}[1]{\vspace{3pt}{\sffamily\bfseries\color{ink}#1}\par\vspace{1pt}}
\newcommand{\vd}[1]{\textsf{\small #1}}

\begin{document}
{\sffamily\color{ink}{\LARGE\bfseries \onus{}}\hfill{\small\href{https://www.onuscheck.com}{www.onuscheck.com} \;·\;
\href{https://github.com/nash911/onus}{github.com/nash911/onus}}\par
\vspace{2pt}{\large Claim-to-evidence maps for research papers, with a measured error rate}\par
\vspace{1pt}{\small Avinash Ranganath \;·\; \href{https://avinashranganath.com}{avinashranganath.com} \;·\;
AI4MetaScience Hackathon 2026, Task 3}\par}
\vspace{2pt}{\color{ink}\rule{\linewidth}{0.8pt}}

\begin{multicols}{2}
\hd{The problem}
Reviewers rarely trace each claim to the evidence behind it, so confident wording passes for support. A
method ``consistently outperforms'' on two of five datasets; an abstract number differs from its table; a
checklist says ``Yes'' to error bars that are not there; code is promised but missing. Fluent AI-written
papers make this cheaper to produce and harder to spot.

\hd{What \onus{} returns}
Give it an arXiv or OpenReview id, or a PDF (optionally with code). It returns a \textbf{claim map}: every
major claim quoted verbatim and located, linked to the table cells, passages, checklist answers and code
lines that bear on it, and given one verdict with a reason: \vd{supported}, \vd{overstated},
\vd{untraced}, \vd{contradicted} or \vd{not\_checkable}. A ledger states what was \emph{not} checked and why.
The report shows the PDF beside the claims; clicking a claim or a piece of evidence highlights it in the paper.
There is no overall score.

\hd{How it works}
\begin{itemize}
  \item \textbf{Code computes, the model points.} Language models extract claims and choose evidence from
    verified candidates. Every quote is checked against the text and every evidence id must come from the
    candidate list. Arithmetic, table comparisons, checklist rules and verdict precedence are plain code.
  \item \textbf{Code checks.} If code is linked or uploaded, \onus{} matches the paper's implementation facts
    to code lines and re-runs the code in an isolated, network-off sandbox.
  \item \textbf{Cheap and fast.} About \texteuro0.10 and 5--8 minutes for a 15-page paper.
\end{itemize}

\hd{How we measured it}
247 Agents4Science 2025 papers, split into dev (10), eval (20) and two test splits (30 each). All
development and tuning used only dev data. The eval split was the pre-freeze go/no-go; the version
was then frozen and fingerprinted, and each test split was run once, three runs per paper. Recall uses
errors planted in held-out papers (five types), against a single-prompt baseline given the same claims.
The protocol and every run are logged publicly in the repository.

\end{multicols}

\vfill
\hd{Results on held-out papers}
\begin{center}
\begin{tabular}{@{}lcccc@{}}
\toprule
 & \multicolumn{2}{c}{test\_1 (frozen, run once)} & \multicolumn{2}{c}{test\_2 (frozen, run once)} \\
\cmidrule(lr){2-3}\cmidrule(lr){4-5}
 & \onus{} & Baseline & \onus{} & Baseline \\
\midrule
\textbf{Claims flagged, unmutated papers} & \textbf{15\%} [12--18] & 29\% [25--33] & \textbf{15\%} [13--18] & 32\% [29--36] \\
Planted-error recall & 70\% [58--80] & 76\% [64--85] & 75\% [62--84] & 78\% [66--87] \\
Collateral flags & 2\% [1--3] & 6\% [5--8] & 3\% [2--4] & 4\% [3--6] \\
Precision, blind labels & 67\% [54--78] & 89\% [72--96] & --- & --- \\
Same verdict in all 3 runs & 84\% & --- & 84\% & --- \\
\bottomrule
\end{tabular}\\[3pt]
{\footnotesize 30 papers per split; Wilson 95\% intervals. Pre-freeze eval split (20 papers): recall 83\% vs.\
80\%, flags 21\% vs.\ 31\%.}
\end{center}

\vspace{-2pt}\noindent\fbox{\parbox{\dimexpr\linewidth-2\fboxsep-2\fboxrule}{\textbf{Half the flags, similar
recall.} On every split \onus{} flags about half as many claims as the baseline, with recall inside the
baseline's interval. For a reviewer, that is attention spent on fewer claims.}}

\vspace{3pt}Its flags are not yet more precise: \vd{untraced} is right only 18 of 31 times, and it misses more
changed abstract numbers (4 of 11 vs.\ 10 of 11 on test\_1). Those are the next things to fix.

\vfill
\begin{multicols}{2}
\hd{Across 90 Agents4Science papers}
\begin{itemize}
  \item Checklist ``Yes'' answers the paper does not back: theory 12 of 18, broader impacts 35 of 69,
    compute 18 of 56, limitations 19 of 81.
  \item 68 papers have code; re-runs were attempted on 48 and only 3 reproduced the paper's values.
  \item A paper's flagged share is unrelated to its review scores ($\rho=-0.05$ AI, $0.10$ human) and to
    acceptance ($p=0.82$).
\end{itemize}
\columnbreak
\hd{Limits}
About 60 planted errors per split (about $\pm$11 points). Planted errors are text-level, five types.
Precision labels were AI-drafted and reviewed by one author. Re-runs are CPU-only and time-capped; live
analyses run once.

\hd{Try it}
\href{https://www.onuscheck.com}{www.onuscheck.com} \;·\; code, protocol and run log:
\href{https://github.com/nash911/onus}{github.com/nash911/onus}
\end{multicols}
\end{document}
